\section{Method Details and End-to-End Comparison}\label{sec:supp-solution}

Section~\ref{sec:method} of the main text gives the overall structure of
candidate-order generation and address assignment as algorithmic pseudocode.
This section expands on three aspects of the design: the rationale for the
complementary candidate orders (\S\ref{sec:supp-framework}), the end-to-end
results compared against the external reference solution
(\S\ref{sec:supp-p1}--\S\ref{sec:supp-headline}), and a diagnostic analysis of
where the gaps come from (\S\ref{sec:supp-failure}).

\subsection{Candidate-Order Design and the Spill-Separation Principle}\label{sec:supp-framework}

The main text treats schedule order as the primary degree of freedom for
optimizing spill cost. This strategy requires the candidate-order set to cover
the different clean/dirty compositions that can arise within peak windows,
rather than merely chasing the local optimum of a single metric. We first
explain why the three candidate orders are complementary, and then discuss the
separation of the prefetch window from spill selection.

\paragraph{Candidate orders.}
All three candidate orders are structure-agnostic topological orders. They share
the same memory-aware list scheduling skeleton (frees-first, allocations-delayed)
and differ only in how they rank the ready set.

\begin{enumerate}
  \item The pressure-aware order (\ksName{p1}) is the memory-aware list
  scheduling order used by Algorithm~\ref{alg:assign-spill} of the main text,
  applying a memory-aware priority within each ready set. \ksOp{alloc} nodes are
  ranked by successor indegree / successor readiness: a lower indegree means the
  consumer becomes ready sooner, so the allocation is deferred until close to its
  use, pushing the allocation timing to the point least sensitive to the peak.
  This order is the direct output of P1 and is also one candidate for P2/P3.
  \item The capacity-throttled order (\ksName{capfit-id}) retains the
  memory-aware free/operation priority within each ready set, but ranks
  \ksOp{alloc} nodes by minimum node identifier instead, and overlays capacity
  gating on top of this base order: if an \ksOp{alloc} would push its cache over
  capacity, the allocation is held back; when all ready allocations are blocked
  and no further deferral is possible, the one that unblocks subsequent work the
  fastest is forced through (with weights normalized by $1/\Capc_c$). The role of
  this order is to directly bound the peak height, at the cost of sacrificing
  some of the continuous residency of buffers.
  \item The ID-reserve order (\ksName{id-raw}) always takes the minimum identifier
  within the free, operation, and allocation ready sets, with neither a
  memory-aware priority nor capacity gating. Its purpose is to preserve the
  original residency order of clean \ksOp{copy-in} buffers, making these clean
  buffers more likely to stay resident through the compute window and thereby
  leaving a clean reserve in high-pressure windows that can be evicted at half
  price. This effect cannot be reliably produced by the first two orders, since
  both of their ranking strategies disturb the original order of the input
  buffers.
\end{enumerate}

The complementarity of the three orders manifests as distinct paths for shaping
the clean/dirty composition: the capacity-throttled order (\ksName{capfit-id})
lowers the total peak through capacity gating, the pressure-aware order
(\ksName{p1}) shortens active intervals through delayed allocation, and the
ID-reserve order (\ksName{id-raw}) maintains a clean reserve through order
preservation. As noted in Section~\ref{sec:liveness-shaping} of the main text,
candidate-set expansion is monotone: adding a new candidate can only improve or
preserve the selected objective.

\paragraph{Address assignment and cost-aware eviction.}
Given a topological order, best-fit address assignment is performed along the
order. When capacity is insufficient, the victim is chosen by a cost-weighted
Belady score:
\begin{equation}\label{eq:supp-victim}
\mathrm{victim}=\arg\max_b d(b)\frac{s_b}{e_b}
=\arg\max_b
\begin{cases}
d(b), & b\ \text{clean},\\
d(b)/2, & b\ \text{dirty},
\end{cases}
\end{equation}
where $d(b)$ is the next-use distance of buffer $b$. The effective distance of a
dirty buffer is halved, so the engine prefers to evict the cheaper clean
buffers: a dirty buffer is evicted ahead of a clean buffer of the same size only
when its future use is far enough away. Proposition~\ref{prop:belady-margin} of
the main text shows that under distance-dominant conditions this scoring rule
selects the same eviction set as other cost-weighted variants, so the marginal
benefit of the eviction policy is far smaller than that of schedule order.

\paragraph{Separating the prefetch window from spill selection.}
The selection of the spill set and the placement of reloads are two separable
degrees of freedom. The former determines the total $\extra$; the latter affects
whether reloads can overlap with computation, and hence the P3 time. Coupling
the two would let time optimization contaminate the traffic objective without
constraint. The method's design choice is to keep them apart. The prefetch
window $H$ controls only the advancement of \ksOp{spill-in}: if an evicted buffer
will be used within $\le H$ steps and the cache has free space, it is reloaded
early, but no extra eviction is triggered to enable the prefetch. P2 uses the
tight grid $H\in\{0,5,40\}$ to minimize traffic; P3 uses the wider grid
$H\in\{0,5,40,80,120\}$ to minimize time.

\paragraph{Selection mechanism.}
A full simulation is run over each point of the Cartesian product of the 3
candidate orders and the prefetch windows: P2 selects by the lexicographic key
$(\extra,\#\mathrm{spills},T)$, and P3 by $(T,\extra,\#\mathrm{spills})$. The
final output is the minimal item over the candidate set.

\begin{algorithm}[htbp]
\caption{Unified solver framework for P1/P2/P3}\label{alg:supp-portfolio}
\begin{algorithmic}[1]
\Require DAG $G=(V,E)$, problem type $p\in\{1,2,3\}$
\Ensure schedule; for P2/P3 also output address assignment $\mathit{mem}$ and spill list $\mathit{spill}$
\If{$p=1$}
  \State \Return $\mathrm{MemoryAwareOrder}(G)$
\EndIf
\State $\mathcal{H}\gets\{0,5,40\}$ \textbf{if} $p=2$ \textbf{else} $\{0,5,40,80,120\}$
\State $\mathit{best}\gets\bot$
\For{$\sigma\in\{\ksName{capfit-id},\ksName{p1},\ksName{id-raw}\}$}
  \For{$H\in\mathcal{H}$}
    \State $(\pi,\mathit{mem},\mathit{spill})\gets \mathrm{AssignAndSpill}(G,\sigma,H)$
    \State $E\gets\sum_{(b,\cdot)\in\mathit{spill}} e_b$;\quad $T\gets \mathrm{SimTime}(\pi,\mathit{mem},\mathit{spill})$
    \State $k\gets(E,|\mathit{spill}|,T)$ \textbf{if} $p=2$ \textbf{else} $(T,E,|\mathit{spill}|)$
    \State $\mathit{best}\gets\min\bigl(\mathit{best},(k,\pi,\mathit{mem},\mathit{spill})\bigr)$
  \EndFor
\EndFor
\State \Return $\mathit{best}$
\end{algorithmic}
\end{algorithm}

\subsection{P1: Peak Residency and Objective Misalignment}\label{sec:supp-p1}

P1 is judged on per-cache peak residency; its result directly answers ``whether
a given DAG can lower its peak under some legal order.'' However, a P1-optimal
order does not automatically translate into a P2/P3 optimum, which is precisely
the objective misalignment phenomenon emphasized in the main text.

The end-to-end results show that our order lowers the L1 peak below the external
reference solution on all six instances: Conv\_Case0 drops from $7488$ to
$6912$, and Conv\_Case1 from $14040$ to $13786$; the reductions are even larger
on the FlashAttention and Matmul instances, where the P1 peaks fall to
$256/256/128/128$. But a low P1 peak is not equivalent to a low spill cost. Take
FlashAttention\_Case1 as an example: the P1 order drives the UB peak to $33010$,
far exceeding $\Capc_{\mathrm{UB}}=1024$. The order that lowers the L1 peak
instead causes a more severe capacity overrun on UB. This case directly explains
why the method must move from the single-peak P1 metric to the spill-traffic
metric of P2: in a multi-cache system, it is the cross-cache distribution of
peaks that determines the true spill cost.

\subsection{P2: Schedule Order Shapes the Clean/Dirty Composition}\label{sec:supp-p2}

The core objective of P2 is to minimize the total spill traffic $\extra$. Unlike
swapping only the eviction policy on a fixed order, our approach uses schedule
order to change the composition of clean and dirty buffers in high-pressure
windows, so that when capacity is exceeded the evictor can free space at a cost
of $s_b$ (clean) rather than $2s_b$ (dirty).

Relative to the external reference solution, P2 wins three and loses three of the
six instances. The wins reveal the effectiveness of the shaping mechanism: in
Conv\_Case1 ($73240\to72520$), our order keeps more clean buffers in the
high-pressure window; Matmul\_Case0 ($34944\to34688$) behaves likewise; and on
Matmul\_Case1 both the $\extra$ and the spill count are level with the reference
solution ($460800$, $3600$), but the third key, time, is better. The three
losing instances point to the flip side of the same mechanism: in the critical
high-pressure window, the reference order happens to retain more suitable clean
buffers. For instance, in Conv\_Case0 the reference solution's $\extra$ is
$73500$, below our $88044$; the gap does not come from a difference in the
eviction rule (Proposition~\ref{prop:belady-margin} shows that different
eviction variants on a fixed order differ very little in extra traffic), but
rather from the reference order keeping more clean buffers resident when spilling
occurs---exactly the shaping effect of schedule order on the residency
composition.

The wins and losses together corroborate the central thesis of the main text:
the primary determinant of spill cost is how the order shapes the clean/dirty
composition, not the eviction rule itself.

\subsection{P3: Prefetch Window and Compute Overlap}\label{sec:supp-p3}

P3 takes pipelined execution time as its primary objective. A nonzero prefetch
window advances some \ksOp{spill-in} operations ahead of their immediate use
points, allowing reloads to overlap with computation or other pipeline transfers
and reducing the wait-blocking of consumer nodes. Because prefetching occurs only
when idle intervals exist and triggers no extra eviction, it mainly adjusts the
timing of reload placement rather than redefining the spill set.

The end-to-end results are four wins and two losses. Our method achieves shorter
times on the two FlashAttention and two Matmul instances: FlashAttention\_Case1
drops from $193059$ to $180364$, and Matmul\_Case0 from $194331$ to $186820$.
Our times lag on the two Conv instances, and Conv\_Case1 is especially
diagnostic: our $\extra$ is lower ($73578<86690$), yet the reference solution's
time is shorter ($1073322<1111778$). This counterexample shows that lower traffic
does not automatically translate into shorter time: the degree of overlap between
reloads and computation is an independent dimension. When the reference order's
reloads happen to be scheduled into idle gaps in the compute pipeline, the
blocking time can be shorter even when the total traffic is higher. This
observation further supports the method's design decision to separate spill
selection from reload placement.

\subsection{End-to-End Reference Comparison}\label{sec:supp-headline}

Table~\ref{tab:supp-headline} summarizes the end-to-end results of the six public
instances on P1/P2/P3, for a combined 13/18 wins. The comparison protocol here
differs from that of the main text's primary experiments: Section~\ref{sec:exp-baselines}
of the main text uses the ``same-engine reorder'' protocol, in which all methods
share the same address assignment / spill engine and only the schedule order is
swapped, so as to isolate the shaping effect of order on the clean/dirty
composition. The end-to-end comparison, by contrast, uses each system's own
complete schedule, address assignment, and spill list, reflecting the final
output quality of two complete systems. The two protocols answer different
questions, and the end-to-end data serve as an overall reference for system
performance.

\begin{table}[H]
\centering
\caption{End-to-end comparison of our method against the external reference solution baseline01 (each cell is ours/reference).}
\label{tab:supp-headline}
\begin{tabular*}{\linewidth}{@{\extracolsep{\fill}}llrrrrc@{}}
\toprule
Case & P & \(\max_{\mathrm{L1}}\) & spills & \(\extra\) & \(T\) & Result \\
\midrule
Conv\_Case0   & P1 & 6912/7488   & 0/0       & 0/0           & 394430/359570   & Win \\
              & P2 & 6912/7488   & 339/141   & 88044/73500   & 615694/535312   & Loss \\
              & P3 & 6912/7488   & 342/141   & 92652/73500   & 553280/535312   & Loss \\
\addlinespace
Conv\_Case1   & P1 & 13786/14040 & 0/0       & 0/0           & 650044/573395   & Win \\
              & P2 & 14040/14040 & 1963/1828 & 72520/73240   & 1159169/1150333 & Win \\
              & P3 & 14040/14040 & 1951/1290 & 73578/86690   & 1111778/1073322 & Loss \\
\addlinespace
FA\_Case0     & P1 & 256/3328    & 0/0       & 0/0           & 31063/31429     & Win \\
              & P2 & 3328/3328   & 31/41     & 3904/3692     & 51700/53374     & Loss \\
              & P3 & 4096/3328   & 58/24     & 5832/4128     & 46167/46761     & Win \\
\addlinespace
FA\_Case1     & P1 & 256/5376    & 0/0       & 0/0           & 112458/113445   & Win \\
              & P2 & 5376/5376   & 229/146   & 32920/32840   & 214361/193059   & Loss \\
              & P3 & 5376/5376   & 229/146   & 34432/32840   & 180364/193059   & Win \\
\addlinespace
Matmul\_Case0 & P1 & 128/9216    & 0/0       & 0/0           & 82308/82773     & Win \\
              & P2 & 9216/9216   & 271/273   & 34688/34944   & 192380/194331   & Win \\
              & P3 & 9216/9216   & 273/273   & 34944/34944   & 186820/194331   & Win \\
\addlinespace
Matmul\_Case1 & P1 & 128/34816   & 0/0       & 0/0           & 629076/583525   & Win \\
              & P2 & 34816/34816 & 3600/3600 & 460800/460800 & 1783145/1800218 & Win \\
              & P3 & 34816/34816 & 3601/3600 & 460928/460800 & 1771132/1800218 & Win \\
\bottomrule
\end{tabular*}
\end{table}

\subsection{Analysis of Where the Gaps Come From}\label{sec:supp-failure}

The core finding across the five losing entries is consistent: the reference
solution's advantage on these instances does not weaken the main text's thesis
but instead corroborates the same mechanism from the opposite direction, namely
that the gap comes mainly from how schedule order shapes the clean/dirty
residency composition, not from a difference in the eviction policy on a fixed
order.

Two kinds of diagnostics cross-validate this judgment. First, when the reference
solution's order is fed into our address assignment / spill engine, the traffic
of several losing entries can be closely reproduced, showing that the gap comes
from the composition of evictable buffers in the high-pressure window rather than
from a difference in the engine implementation. Second, swapping among four
Belady-style victim policies on a fixed order produces only small differences in
P2 $\extra$ (Section~\ref{sec:exp-mechanism} of the main text reports
CV $\le 4\%$); by contrast, $\extra$ can differ by an order of magnitude across
different legal orders.

We keep these losing entries rather than fitting them with operator-family
specialization rules, because such fitting would break the method's
structure-agnostic property: the algorithm uses only the DAG dependencies, buffer
sizes, cache types, clean/dirty attributes, and capacity information, and
contains no ranking rules tailored to operator categories. Once specialization is
introduced, the conclusion would degrade from ``schedule order shapes the
clean/dirty composition to lower spill cost'' to ``a dedicated ordering for a
specific operator structure,'' weakening the generality of the method.

\subsection{Runtime Efficiency}\label{sec:supp-runtime}

Table~\ref{tab:supp-runtime} gives the complete statistics over 3 repetitions
(median, min, max), supplementing the median data of Table~\ref{tab:runtime} in
the main text.

\begin{table}[htbp]
\centering\small
\caption{Wall-clock runtime (seconds, 3 repetitions).}\label{tab:supp-runtime}
\begin{tabular*}{\linewidth}{@{\extracolsep{\fill}}llrrrr@{}}
\toprule
Instance & P & Median & Min & Max & Nodes \\
\midrule
Conv\_Case0 & P1 & 0.009 & 0.009 & 0.010 & 2\,580 \\
            & P2 & 0.391 & 0.390 & 0.392 & \\
            & P3 & 0.683 & 0.677 & 0.706 & \\
\addlinespace
Conv\_Case1 & P1 & 0.189 & 0.186 & 0.190 & 36\,086 \\
            & P2 & 57.97 & 57.46 & 58.62 & \\
            & P3 & 73.14 & 73.10 & 73.28 & \\
\addlinespace
FA\_Case0   & P1 & 0.006 & 0.006 & 0.007 & 1\,716 \\
            & P2 & 0.220 & 0.218 & 0.222 & \\
            & P3 & 0.328 & 0.326 & 0.328 & \\
\addlinespace
FA\_Case1   & P1 & 0.027 & 0.026 & 0.030 & 6\,952 \\
            & P2 & 2.723 & 2.723 & 2.757 & \\
            & P3 & 3.356 & 3.316 & 3.367 & \\
\addlinespace
Matmul\_Case0 & P1 & 0.014 & 0.014 & 0.016 & 4\,160 \\
              & P2 & 0.701 & 0.695 & 0.703 & \\
              & P3 & 1.137 & 1.136 & 1.170 & \\
\addlinespace
Matmul\_Case1 & P1 & 0.126 & 0.111 & 0.126 & 30\,976 \\
              & P2 & 19.40 & 19.33 & 19.46 & \\
              & P3 & 32.33 & 32.19 & 32.34 & \\
\bottomrule
\end{tabular*}
\end{table}

\paragraph{Dominant factors in runtime.}
P1 generates only a single schedule order, and all instances complete in under
$0.2$ seconds. The runtime of P2/P3 is dominated by two factors: (1) the grid
size of candidate orders $\times$ prefetch windows (P2 evaluates $3\times3=9$
combinations, P3 evaluates $3\times5=15$ combinations); and (2) the quadratic
worst-case cost of Belady eviction selection in address assignment. Conv\_Case1
and Matmul\_Case1 have 12\,013 and 8\,960 buffers respectively, and the eviction
scan constitutes the bulk of their runtime. The min/max values over the three
repetitions tightly bracket the median (coefficient of variation $<2\%$),
indicating that the timing results are stable.

\paragraph{Time complexity.}
Let $N=|V|$ be the number of nodes, $E=|E|$ the number of edges, $B$ the number
of buffers, $K=3$ the number of candidate orders, and $W$ the size of the
prefetch-window grid. The complexity of each stage is as follows:
\begin{itemize}
  \item \textbf{Candidate-order generation}: $O(K(N\log N+E))$, a heap-maintained
  ready-set topological sort.
  \item \textbf{Address assignment and spill insertion}: $O(N+E+B^2)$ per
  $(\text{order},\text{window})$ combination. The sweep along the order is linear,
  and the worst-case cost of Belady eviction is $B$ allocations each scanning
  $O(B)$ resident buffers.
  \item \textbf{Pipeline simulation}: $O(N+E)$ per combination.
  \item \textbf{Portfolio selection}: enumeration of $O(KW)$ combinations.
\end{itemize}
The total is $O(KW(N\log N+E+B^2))$. Since $K$ and $W$ are small constants, the
dominant asymptotic term is $O(N\log N+E+B^2)$. The space complexity is
$O(N+E+B)$.
